% Generated from locale/tr/content/first-order-logic/axiomatic-deduction/proving-things.tex; do not hand-edit.


\iftag{FOL}
      {\olfileid[tr]{fol}{axd}{pro}}
      {\olfileid[tr]{pl}{axd}{pro}}

\olsection{\usetoken{P}{derivation} Örnekleri}

\begin{ex}
$(\lnot !D \lor !E) \lif (!D \lif !E)$ ifadesini ispatlamak istediğimizi
varsayalım. Bunun aksiyomlarımızdan herhangi birinin örneği olmadığı açıktır;
dolayısıyla onu !!{derive} işlemiyle elde etmek için \MP{} kuralını
kullanmalıyız. Bu örnekte gereken kural MP'dir; $!A$ ile $!A \lif !B$
verildiğinde $!B$'yi doğru bir adım olarak gerekçelendirmemizi sağlar. Bir strateji,
\olref[prp]{ax:lor3} içinde $!A$ yerine $\lnot !D$, $!B$ yerine $!E$ ve $!C$
yerine $!D \lif !E$ koymak, yani şu örneği kullanmak olabilir:
\[
(\lnot !D \lif (!D \lif !E)) \lif ((!E \lif (!D \lif !E)) \lif ((\lnot
!D \lor !E) \lif (!D \lif !E))).
\]
Neden? MP'nin iki kez uygulanması, istediğimiz son bölümü verir.
$\lnot !D \lif (!D \lif !E)$ ifadesinin \olref[prp]{ax:lnot2}'nin ve
$!E \lif (!D \lif !E)$ ifadesinin \olref[prp]{ax:lif1}'in bir örneği
olduğunu kolayca görürüz. Böylece !!{derivation} şöyledir:
\begin{derivation}
  1. &  $\lnot !D \lif (!D \lif !E)$ & \olref[prp]{ax:lnot2} \\
  2. & $(\lnot !D \lif (!D \lif !E)) \lif {}$\\
  &\qquad $((!E \lif (!D \lif !E)) \lif ((\lnot !D \lor !E) \lif (!D \lif !E)))$ & \olref[prp]{ax:lor3}\\
  3. & $(!E \lif (!D \lif !E)) \lif ((\lnot !D \lor !E) \lif (!D \lif !E))$ & 1, 2, \MP\\
  4. & $!E \lif (!D \lif !E)$ & \olref[prp]{ax:lif1}\\
  5. & $(\lnot !D \lor !E) \lif (!D \lif !E)$ & 3, 4, \MP
\end{derivation}
\end{ex}

\begin{ex}\ollabel{ex:identity}
$!D \lif !D$ için bir !!a{derivation} bulmaya çalışalım. Bu, bir aksiyom
örneği değildir; dolayısıyla onu !!{derive} işlemiyle elde etmek için \MP{}
kullanmalıyız. \olref[prp]{ax:lif1}, \MP{} uygulayabileceğimiz $!A \lif !B$
biçiminde bir aksiyomdur. Bunun işe yaraması için \MP{}'nin doğru bir adım
olarak vereceği $!B$ elbette hedefimiz olan $!D \lif !D$ olmalıdır. O zaman
$!A$ da $!D$ olmalı; yani \olref[prp]{ax:lif1}'in şu örneğine bakabiliriz:
\[
!D \lif (!D \lif !D)
\]
\MP{} uygulamak için karşılık gelen ikinci öncülü, yani $!A$'yı da
gerekçelendirmemiz gerekir. Bizim durumumuzda bu $!D$ olur; ancak
$!D$'yi tek başına !!{derive} işlemiyle elde edemeyiz. Bu nedenle başka bir
strateji gerekir.

Yalnızca $\lif$ içeren diğer aksiyom \olref[prp]{ax:lif2}, yani
\[
(!A \lif (!B \lif !C)) \lif ((!A \lif !B) \lif (!A \lif !C))
\]
ifadesidir. İç içe son koşula \MP{}'yi iki kez uygulayarak ulaşabiliriz.
Dolayısıyla \olref[prp]{ax:lif2}'nin, $!A \lif !C$ bölümünün !!{derive}
işlemiyle elde etmek istediğimiz !!{formula} olan $!D \lif !D$'ye eşit
olduğu bir örneğini isteriz. Elbette
$!A$ ile $!C$ bu durumda $!D$'dir. $!A \lif (!B \lif !C)$ ile
$!A \lif !B$ ifadelerinin, yani $!D \lif (!B \lif !D)$ ile
$!D \lif !B$ ifadelerinin de !!{derivable} olması için $!B$'yi nasıl
seçmeliyiz? Bunların ilki, $!B$ için ne seçersek seçelim zaten
\olref[prp]{ax:lif1}'in bir örneğidir. $!B$'yi $(!D \lif !D)$ seçersek
$!D \lif !B$ de \olref[prp]{ax:lif1}'in başka bir örneği olur. Buna göre
!!{derivation} şöyledir:
\begin{derivation}
  1. & $!D \lif ((!D \lif !D) \lif !D)$ & \olref[prp]{ax:lif1}\\
  2. & $(!D \lif ((!D \lif !D) \lif !D)) \lif {}$\\
  & \qquad $((!D \lif (!D \lif !D)) \lif (!D \lif !D))$ & \olref[prp]{ax:lif2}\\
  3. & $(!D \lif (!D \lif !D)) \lif (!D \lif !D)$ & 1, 2, \MP\\
  4. & $!D \lif (!D \lif !D)$ & \olref[prp]{ax:lif1}\\
  5. & $!D \lif !D$ & 3, 4, \MP
\end{derivation}
\end{ex}

\begin{ex}\ollabel{ex:chain}
Bazen bir !!a{derivation} bulunduğunu göstermek isteriz: hedef bir
!!{formula}, $\Gamma$ ise başka !!{formula}s içerir. Örneğin
$\Gamma = \{!A \lif !B, !B \lif !C\}$ kümesinden $!A \lif !C$'yi
!!{derive} işlemiyle elde edebileceğimizi gösterelim.
\begin{derivation}
  1. & $!A \lif !B$ & \Hyp\\
  2. & $!B \lif !C$ & \Hyp\\
  3. & $(!B \lif !C) \lif (!A \lif (!B \lif !C))$ & \olref[prp]{ax:lif1} \\
  4. & $!A \lif (!B \lif !C)$ & 2, 3, \MP\\
  5. & $(!A \lif (!B \lif !C)) \lif {}$\\
  & \qquad $((!A \lif !B) \lif (!A \lif !C))$ & \olref[prp]{ax:lif2}\\
  6. & $((!A \lif !B) \lif (!A \lif !C))$ & 4, 5, \MP\\
  7. & $!A \lif !C$ & 1, 6, \MP
\end{derivation}
``\Hyp'' (hypothesis, yani varsayım) etiketi taşıyan satırlar, satırdaki
ifadenin bir !!{formula} olduğunu ve $\Gamma$ içinde bir !!a{element} olarak
yer aldığını belirtir.
\end{ex}

\begin{prop}
  \ollabel{prop:chain} $\Gamma \Proves !A \lif !B$ ve $\Gamma
  \Proves !B \lif !C$ ise $\Gamma \Proves !A \lif !C$.
\end{prop}

\begin{proof}
  $\Gamma \Proves !A \lif !B$ ve $\Gamma \Proves !B \lif !C$ olduğunu
  varsayalım. O zaman $\Gamma$'dan $!A \lif !B$ için bir !!a{derivation};
  ayrıca $\Gamma$'dan $!B \lif !C$ için bir !!a{derivation} vardır. Bunları
  art arda yazarak tek bir !!{derivation} içinde birleştirin. Şimdi önceki
  örnekteki !!{derivation} içinden 3--7. satırları ekleyin. Böylece $\Gamma$'dan
  şu !!a{derivation} elde edilir; bunun son satırı $!A \lif !C$'dir. Bu
  !!{derivation} için 4. ve 7. satırların gerekçeleri geçerliliğini
  korur: 1. satıra yapılan başvuruyu $!A \lif !B$ ile biten !!{derivation} içindeki son
  satıra, 2. satıra yapılan başvuruyu da $!B \lif !C$'nin
  ilgili !!{derivation} içindeki son satıra yapılan başvuruyla değiştiririz.
\end{proof}

\begin{prob}
Aşağıdakileri aksiyomlardan başlayan !!{derivation}s sergileyerek gösterin:
\begin{enumerate}
\item $(!A \land !B) \lif (!B \land !A)$
\item $((!A \land !B) \lif !C) \lif (!A \lif (!B \lif !C))$
\item $\lnot(!A \lor !B) \lif \lnot !A$
\end{enumerate}
\end{prob}

